\subsection{Character, trace identities, and the abelian index}

\begin{proposition}[Split weight enumerator and one-family character]
\label{prop:split-weight-enumerator}
For the present paragraph let
\[
\chi_{S^+}(t):=\tr_{S^+}(e^{tY}).
\]
Define also the split weight enumerator
\[
W(u,v)
:=
\sum_{\substack{0\le n_3\le 3\\ 0\le n_2\le 2\\ n_3+n_2\ \mathrm{even}}}
\binom{3}{n_3}\binom{2}{n_2}u^{n_3}v^{n_2}.
\]
Then
\begin{equation}
\label{eq:split-weight-enumerator}
W(u,v)
=
1+3u^2+6uv+v^2+2u^3v+3u^2v^2
=
\frac12\Big[(1+u)^3(1+v)^2+(1-u)^3(1-v)^2\Big].
\end{equation}
Moreover,
\begin{equation}
\label{eq:family-character}
\chi_{S^+}(t)
=
\frac12\Big[
\bigl(1+e^{-t/3}\bigr)^3\bigl(1+e^{t/2}\bigr)^2
+
\bigl(1-e^{-t/3}\bigr)^3\bigl(1-e^{t/2}\bigr)^2
\Big]
=
W(e^{-t/3},e^{t/2}).
\end{equation}
Consequently,
\[
\chi_{S^+}(0)=16,
\qquad
\chi_{S^+}'(0)=\tr_{S^+}Y=0,
\qquad
\chi_{S^+}''(0)=\tr_{S^+}Y^2=\frac{10}{3},
\qquad
\chi_{S^+}'''(0)=\tr_{S^+}Y^3=0.
\]
\end{proposition}

\begin{remark}[Master generating function for the one-family packet]
The even-parity split weight enumerator
\[
W(u,v)=\frac12\bigl[(1+u)^3(1+v)^2+(1-u)^3(1-v)^2\bigr]
\]
is the master generator of the family packet. The one-family character is obtained by the
specialization $u=e^{-t/3}$ and $v=e^{t/2}$, while all trace moments of $Y$ are recovered by
differentiating at $t=0$.
\end{remark}

\begin{proof}
The pairs $(n_3,n_2)$ allowed by even total degree are
\[
(0,0),\ (2,0),\ (1,1),\ (0,2),\ (3,1),\ (2,2),
\]
with multiplicities
\[
1,\ 3,\ 6,\ 1,\ 2,\ 3.
\]
This gives the explicit polynomial formula for $W(u,v)$. Equivalently, the parity projector
\[
\frac12\bigl(1+(-1)^{n_3+n_2}\bigr)
\]
acting on $(1+u)^3(1+v)^2$ yields \Cref{eq:split-weight-enumerator}. Substituting
\[
u=e^{-t/3},
\qquad
v=e^{t/2}
\]
gives \Cref{eq:family-character}. For an even exterior algebra one also has
\[
\tr_{\Lambda^{\mathrm{even}}E}(e^{tY})
=
\frac12\Big[\prod_{j=1}^5(1+e^{tq_j})+\prod_{j=1}^5(1-e^{tq_j})\Big],
\]
with carrier charges $q_j\in\{-1/3,-1/3,-1/3,1/2,1/2\}$. Differentiating at $t=0$ yields the stated dimension and trace identities.
\end{proof}

\begin{corollary}[Multiplicative character as six-sector decoder]
\label{cor:multiplicative-family-character}
Let
\[
X:=6Y.
\]
Then the multiplicative one-family character
\[
\Xi_{S^+}(t):=\tr_{S^+}(t^X)
\]
is
\begin{equation}
\label{eq:multiplicative-family-character}
\Xi_{S^+}(t)
=
\frac12\Big[
(1+t^{-2})^3(1+t^3)^2
+
(1-t^{-2})^3(1-t^3)^2
\Big]
=
t^6+3t^2+6t+1+2t^{-3}+3t^{-4}.
\end{equation}
Consequently
\[
\Xi_{S^+}(1)=16,
\qquad
(t\partial_t)\Xi_{S^+}(1)=\tr_{S^+}X=0,
\qquad
(t\partial_t)^2\Xi_{S^+}(1)=\tr_{S^+}X^2=120=5!.
\]
Thus the six sectors and their multiplicities are read off directly from the monomials of
\Cref{eq:multiplicative-family-character}.
\end{corollary}

\begin{proof}
Since $X=3N_2-2N_3$, every admissible occupancy pair contributes the monomial
$t^{3n_2-2n_3}$. Therefore
\[
\Xi_{S^+}(t)=W(t^{-2},t^3),
\]
and \Cref{eq:split-weight-enumerator} gives the displayed factorized form. Expanding yields the six
monomials shown above. Applying $t\partial_t$ once and twice at $t=1$ recovers the first two trace
moments of $X$.
\end{proof}

\begin{remark}[Table as readout, character as proof device]
The one-family table remains useful as a reader-facing decoding sheet, but the whole six-sector
arithmetic is already compressed into the single multiplicative character
\Cref{eq:multiplicative-family-character}.
\end{remark}

\begin{remark}[Six admissible occupancy sectors]
The six monomials in \Cref{eq:split-weight-enumerator} are exactly the six allowed occupancy pairs $(n_3,n_2)$ of the family packet. This is the shortest combinatorial route from the split carrier to the six Standard Model multiplet sectors listed in \TFPTxref{tab:one-family}.
\end{remark}

The carrier packet therefore reproduces the standard trace cancellations
\[
\tr_{S^+}Y=0,
\qquad
\tr_{S^+}Y^3=0,
\]
and packages them together with the quadratic trace in one formula. Since $\tr_{S^+}Y^2=4\gamma=10/3$ and $\tr_\Phi Y^2=1/2$ for one weak doublet with hypercharge $1/2$, the hypercharge trace coefficient organizes as
\begin{equation}
\label{eq:abelian-index}
\bone(N_{\mathrm{fam}},N_\Phi)
=
\frac25N_{\mathrm{fam}}\tr_{S^+}Y^2+\frac15N_\Phi\tr_\Phi Y^2
=
\frac85N_{\mathrm{fam}}\gamma+\frac1{10}N_\Phi
=
\frac43N_{\mathrm{fam}}+\frac1{10}N_\Phi.
\end{equation}
On the canonical branch $(N_{\mathrm{fam}},N_\Phi)=(3,1)$ this gives
\[
\bone=\frac{41}{10}.
\]
The abelian index is therefore not inserted as an isolated Standard Model datum; it is a carrier-level output once the family and Higgs counts are declared.

\begin{remark}[Canonical-branch Diophantine lock]
If one imposes $\bone=41/10$ on nonnegative integer data, then
\[
40N_{\mathrm{fam}}+3N_\Phi=123.
\]
The physically nontrivial solution is $(N_{\mathrm{fam}},N_\Phi)=(3,1)$, so the canonical branch is already sharply constrained at the level of the trace arithmetic.
\end{remark}

\begin{proposition}[Family hypercharge arithmetic and compression]
\label{prop:family-hypercharge-compression}
Let $N_3$ and $N_2$ be the occupation-number operators on the $E_3$ and $E_2$ slots, restricted to $S^+$, and define the scaled hypercharge operator
\[
X:=6Y=3N_2-2N_3.
\]
Then on $S^+$ one has
\begin{align}
X&\equiv N_2\equiv N_3 \pmod 2,\label{eq:x-mod-2}\\
X&\equiv N_3 \pmod 3,\label{eq:x-mod-3}\\
X&\equiv 3(N_2+N_3) \pmod 5.\label{eq:x-mod-5}
\end{align}
In particular, the even exterior degree $k=N_2+N_3\in\{0,2,4\}$ is read off from the residue class of $X$ modulo $5$: if $X\equiv r\pmod 5$ with $r\in\{0,1,2\}$, then $k=2r$, so
\[
X\bmod 5=0,1,2
\qquad\Longleftrightarrow\qquad
\Lambda^0E,\ \Lambda^2E,\ \Lambda^4E.
\]
Moreover,
\[
\mathrm{Spec}(X|_{S^+})=\{0,-4,1,6,-3,2\},
\]
and therefore
\begin{equation}
\label{eq:family-x-polynomial}
X(X+4)(X-1)(X-6)(X+3)(X-2)=0
\qquad\text{on }S^+.
\end{equation}
\end{proposition}

\begin{proof}
Each occupied $E_3$ mode contributes hypercharge $-1/3$ and each occupied $E_2$ mode contributes $+1/2$, so multiplying by $6$ gives $X=3N_2-2N_3$. Modulo $2$, the term $-2N_3$ vanishes, so $X\equiv N_2$. Because $S^+$ has even total degree, $N_2\equiv N_3\pmod 2$, which proves \Cref{eq:x-mod-2}. Modulo $3$, one has $3N_2\equiv 0$ and $-2\equiv 1$, giving \Cref{eq:x-mod-3}. Modulo $5$, the congruence $-2\equiv 3$ yields \Cref{eq:x-mod-5}. Since on $S^+$ the even degree is $k=N_2+N_3\in\{0,2,4\}$, the possible residues are $3k\equiv 0,1,2\pmod 5$, from which $k=2r$ with $r=X\bmod 5$ follows. Evaluating $X=3n_2-2n_3$ on the six admissible occupancy pairs $(n_3,n_2)$ gives the listed spectrum. The polynomial \Cref{eq:family-x-polynomial} is then the spectral polynomial with these six simple roots.
\end{proof}

\begin{corollary}[Five-fold phase filters for even degree]
\label{cor:family-degree-filters}
Let
\[
\zeta_5=e^{2\pi i/5},
\qquad
U_5:=e^{2\pi iX/5}.
\]
Then the even-degree projectors are the discrete Fourier filters
\[
P_{\Lambda^{2m}E}
=
\frac15\sum_{r=0}^{4}\zeta_5^{-mr}U_5^r,
\qquad
m=0,1,2.
\]
\end{corollary}

\begin{proof}
By \TFPTxref{prop:family-hypercharge-compression}, the eigenspaces $\Lambda^{2m}E$ carry the residue class $X\equiv m\pmod 5$. Hence $U_5$ has eigenvalue $\zeta_5^m$ on $\Lambda^{2m}E$, and the standard Fourier projector formula on the cyclic group $\mathbb Z_5$ gives the claim.
\end{proof}

\begin{remark}[From the two-point carrier to the six-point family algebra]
The minimal carrier itself is a two-point algebra, but the family packet lifts the scaled hypercharge to a six-point algebra. Therefore every analytic function of $X$ on $S^+$ reduces to a polynomial of degree at most five. If
\[
M_n:=\tr_{S^+}(X^n),
\]
then \Cref{eq:family-x-polynomial} implies the finite recursion
\[
M_{n+6}
=
2M_{n+5}+31M_{n+4}-20M_{n+3}-156M_{n+2}+144M_{n+1}.
\]
In this sense the whole hypercharge statistics of one family is finitely compressed into six spectral values.
\end{remark}


% ---- Compression identities and carrier-relevant comparison setup: source lines 2696-3156 ----
\section{Compression identities and numerical comparison of alternative discrete worlds}

This section collects the algebraic compressions that turn the carrier packet and kernel numerics
from a list of internally consistent outputs into an overdetermined structure. These compression
identities are derived consistency statements on the retained branch: their formal derivation closes
only once the bosonic-rank and branch-Yukawa carrier theorems have fixed the rigid $3+2$ split. By
contrast, the discrete-world table below is a numerical comparison layer: it is not part of the
rigid theorem surface and should be read as an appendix-style pressure test on the carrier packet.

